% !TEX root = main.tex
\chapter{Συμπεράσματα και Συζήτηση}
\label{chap:conclusions}

Το κεφάλαιο αυτό συγκεντρώνει τα κύρια ευρήματα του
Κεφαλαίου~\ref{chap:experiments}, ελέγχει την υπόθεση της
Ενότητας~\ref{sec:intro-hypothesis}, δίνει μια σύσταση ανά πηγή δυσκολίας και
παραθέτει πιθανές κατευθύνσεις για μελλοντική έρευνα.

\section{Επανεξέταση της Κεντρικής Υπόθεσης}
\label{sec:concl-hypothesis}

Η υπόθεση της Ενότητας~\ref{sec:intro-hypothesis} περιείχε επτά ισχυρισμούς. Τους
εξετάζουμε με βάση τα αποτελέσματα των πειραμάτων. Στο τέλος της ενότητας δίνουμε μια
παρατήρηση για τη μεροληψία του ULA που προέκυψε από τα αποτελέσματα.

\paragraph{Ισχυρισμός 1: Ο ULA είναι ασυμπτωτικά μεροληπτικός.}
\emph{Επαληθεύεται.} Στη Γκαουσιανή με $\eta = 0.3$ η διασπορά των δειγμάτων του ULA
σταθεροποιείται στο $1.178$, την τιμή $1/(1-\eta/2)$ που προβλέπει η θεωρία, και
μένει εκεί έως $10^6$ επαναλήψεις (Σχήμα~\ref{fig:case-ula-vs-mala}). Στη μελέτη
σύγκλισης ο λόγος του ULA απέχει πολύ από τον ιδανικό σε τέσσερις κατανομές: στη
Γκαουσιανή η KS πέφτει μόνο από $0.0131$ σε $0.0073$ (λόγος $0.55$, ιδανικό $0.224$),
και στο μείγμα το $b^2$ δεν μειώνεται μετά τις $10^5$ επαναλήψεις (μένει στο $0.0078$). Ο ULA συγκλίνει σε μια
στάσιμη κατανομή που δεν συμπίπτει με την κατανομή-στόχο.

\paragraph{Ισχυρισμός 2: Η διόρθωση Metropolis στον MALA εξαλείφει τη μεροληψία
του ULA.}
\emph{Επαληθεύεται.} Με το ίδιο βήμα και την ίδια πρόταση ο MALA δίνει διασπορά
$1.000$ στη Γκαουσιανή. Στη μελέτη σύγκλισης ο MALA έχει λόγο κοντά στον ιδανικό σε
τέσσερις από τις επτά κατανομές: KS $0.22$ στη Γκαουσιανή και $0.20$ στη Student-$t$,
$b^2$ $0.057$ στη μπανάνα και $0.042$ στο διπλό πηγάδι. Το σφάλμα που απομένει είναι
θόρυβος Monte Carlo. Στο μείγμα το $b^2$ του σταματά στο $0.0011$, και στην Cauchy
κάποιες αλυσίδες επιστρέφουν από τις ουρές μετά τις $5 \cdot 10^4$ επαναλήψεις. Η εξαίρεση είναι το χωνί, όπου όλοι οι
αλγόριθμοι αναμειγνύονται πολύ αργά.

\paragraph{Ισχυρισμός 3: Οι βαριές ουρές δυσκολεύουν τις overdamped παραλλαγές.}
\emph{Επαληθεύεται για τον ULA, εν μέρει για τον MALA.} Στην Cauchy με $\eta = 0.1$ στο
$d=10$ ο ULA καλύπτει μόνο το $2\%$ της σωστής μάζας πέρα από τα ποσοστημόρια $1\%$
και $99\%$, και το σφάλμα ποσοστημορίου στο $97.5\%$ είναι $6.7$, περίπου εκατό φορές
μεγαλύτερο από της Γκαουσιανής. Ο MALA αντισταθμίζει με μεγαλύτερο βήμα (πέντε φορές
μεγαλύτερο από του ULA στο $d=10$) και στο $d=20$ φτάνει KS $0.026$ στις $10^6$
επαναλήψεις.
Για $d \ge 100$ όμως η KS είναι πάνω από $0.2$ για κάθε αλγόριθμο.

\paragraph{Ισχυρισμός 4: Το kinetic σχήμα βελτιώνει τη δειγματοληψία βαριών ουρών.}
\emph{Επαληθεύεται εν μέρει.} Ο BAOAB υπερέχει του ULA στην Cauchy σε κάθε διάσταση.
Έναντι του MALA δεν έχει σταθερό πλεονέκτημα. Στο $d=10$ οι δύο αλγόριθμοι έχουν KS μέσα στο IQR ο ένας του άλλου ($0.034$ και $0.038$). Στο
$d=20$, στις $10^6$ επαναλήψεις, ο BAOAB έχει μικρότερη KS ($0.020$ έναντι $0.026$),
αλλά η KS του δεν μειώνεται μετά τις $5 \cdot 10^5$ επαναλήψεις, ενώ του MALA
μειώνεται ακόμα. Σε υψηλή διάσταση και οι τρεις
αποτυγχάνουν.

\paragraph{Ισχυρισμός 5: Η επιλογή διακριτοποίησης έχει σημασία.}
\emph{Επαληθεύεται.} Ο ULA και ο BAOAB δεν έχουν έλεγχο αποδοχής και χρησιμοποιούν την
ίδια κλίση. Στο $d=10$ ο BAOAB έχει μικρότερη κύρια μετρική από τον ULA σε επτά από τις
οκτώ κατανομές. Η εξαίρεση είναι το διπλό πηγάδι, όπου ο ULA κερδίζει μόνο επειδή η
μεροληψία του χαμηλώνει το φράγμα (βλ. την παρατήρηση στο τέλος της ενότητας). Στη Γκαουσιανή ο BAOAB δεν έχει
μεροληψία στη θέση για κανένα ευσταθές βήμα \citep{leimkuhler2013bao}, και η KS του
είναι στο επίπεδο ανεξάρτητων δειγμάτων.

\paragraph{Ισχυρισμός 6: Η ανθεκτικότητα του BAOAB στη διάσταση ισχύει και χωρίς
ισοτροπία.}
\emph{Επαληθεύεται, με μία εξαίρεση.} Ο BAOAB μένει σταθερός έως $d=1000$ στη
Γκαουσιανή (KS $0.004$), στην ανισοτροπική Γκαουσιανή ($b^2 \approx 0.001$), στο
μείγμα ($b^2 \approx 0.005$) και στο διπλό πηγάδι ($b^2 \approx 0.005$), και στη
μπανάνα έως $d=200$. Η εξαίρεση είναι το χωνί, όπου όλες οι συντεταγμένες εξαρτώνται
από την ίδια $v$ και η δυσκολία μεγαλώνει με τη διάσταση: εκεί ο BAOAB είναι ο
καλύτερος έως $d=500$, αλλά αποτυγχάνει για $d \ge 20$.

\paragraph{Ισχυρισμός 7: Ο MALA χειροτερεύει με τη διάσταση σε σταθερό μήκος
αλυσίδας, χωρίς μεροληψία.}
\emph{Επαληθεύεται.} Σε $5 \cdot 10^4$ επαναλήψεις ο MALA χειροτερεύει με τη διάσταση
σε κάθε κατανομή, και το βέλτιστο βήμα του μικραίνει σε κάθε κατανομή (για παράδειγμα
από $0.5$ σε $0.1$ στη Γκαουσιανή, από $2$ σε $0.05$ στη μπανάνα). Στο $d=20$ όμως ο
MALA είναι τελευταίος στη μπανάνα και στο διπλό πηγάδι στις $5 \cdot 10^4$ επαναλήψεις.
Στις $10^6$ είναι δεύτερος στη μπανάνα και ισόπαλος με τον ULA στο διπλό πηγάδι, με
τον ιδανικό ρυθμό. Η κακή εικόνα του MALA σε μικρές αλυσίδες σημαίνει
ότι χρειάζεται μεγαλύτερη αλυσίδα, όχι ότι δειγματοληπτεί λάθος κατανομή.

\paragraph{Παρατήρηση: η μεροληψία του ULA δεν φαίνεται πάντα στην κύρια μετρική.}
Στην ανισοτροπική Γκαουσιανή (για $d \ge 20$) και στο μείγμα το
$b^2$ του ULA είναι κάτω από το κατώφλι ($0.0095$ και $0.008$), αλλά η κάλυψη ουράς του είναι $1.1$--$1.2$
και $1.6$. Στο διπλό πηγάδι ο ULA έχει το μικρότερο $b^2$ από τους τρεις, γιατί
βάζει $1.79$ φορές τη σωστή μάζα κοντά στο φράγμα και το διασχίζει συχνότερα· η κάλυψη
ουράς του είναι $1.5$--$1.8$, και στις $10^6$ επαναλήψεις χάνει την πρώτη θέση. Ούτε το
ESS ούτε το σφάλμα στη μάζα κάθε κορυφής δείχνουν αυτή τη μεροληψία.

\section{Σύσταση ανά Πηγή Δυσκολίας}
\label{sec:concl-recommendation}

Ο Πίνακας~\ref{tab:recommendation} συνοψίζει τα ευρήματα ανά πηγή δυσκολίας: ποιος
αλγόριθμος είναι ο καλύτερος στο $d=10$, πού αποτυγχάνει κάθε αλγόριθμος, και πώς
επιδρά η διάσταση.

\begin{table}[H]
\centering
\footnotesize
\begin{tabular}{p{2.3cm}p{1.7cm}p{6.0cm}p{3.6cm}}
\toprule
Πηγή δυσκολίας & Καλύτερος & Πού αποτυγχάνει κάθε αλγόριθμος & Επίδραση της διάστασης \\
\midrule
Καμία (Γκαουσιανή) & BAOAB & ULA: μεροληψία στη διασπορά. MALA: δεύτερος σε μικρή αλυσίδα. & BAOAB σταθερός έως $d=1000$· MALA χειροτερεύει αργά. \\
\addlinespace
Μέτριες και βαριές ουρές (Student-$t$, Cauchy) & BAOAB, MALA & ULA: δεν φτάνει στις ουρές της Cauchy σε μικρή διάσταση, βάζει πολλή μάζα στις ουρές της Student-$t$ σε μεγάλη. MALA και BAOAB: ισόπαλοι στην Cauchy στο $d=10$· ο MALA συγκλίνει ταχύτερα. & Στην Cauchy όλοι αποτυγχάνουν για $d \ge 100$. \\
\addlinespace
Διαφορετικές κλίμακες & BAOAB & ULA: αργή κίνηση στη φαρδιά κατεύθυνση, μεροληψία κάτω από το κατώφλι. MALA: αποτυγχάνει για $d \ge 500$. & BAOAB σταθερός έως $d=1000$. \\
\addlinespace
Καμπύλη συσχέτιση & BAOAB & ULA: αποτυγχάνει για $d \ge 4$. MALA: καταρρέει σε υψηλή διάσταση ($b^2 = 2.9$ στο $d=1000$). & BAOAB σταθερός έως $d=200$. \\
\addlinespace
Πολλές κορυφές & BAOAB & ULA: μεροληψία στη διασπορά κάτω από το κατώφλι. MALA: λιγότερες μεταβάσεις ανάμεσα στις κορυφές, αλλά σωστός με μεγαλύτερη αλυσίδα. & BAOAB σταθερός έως $d=1000$. \\
\addlinespace
Κλίμακα που αλλάζει με τη θέση & BAOAB & ULA: πολύ μικρό ευσταθές βήμα, δεν φτάνει στο στόμιο. Όλοι: πολύ αργή ανάμειξη, και στις $10^6$ επαναλήψεις. & Όλοι αποτυγχάνουν για $d \ge 20$. \\
\addlinespace
Φράγμα, μη λείο δυναμικό & BAOAB & ULA: «κερδίζει» το $b^2$ μόνο λόγω μεροληψίας στη μάζα του φράγματος. MALA: σπάνιες διασχίσεις, σωστός με μεγαλύτερη αλυσίδα. & BAOAB σταθερός έως $d=1000$· MALA χειροτερεύει. \\
\bottomrule
\end{tabular}
\caption{Σύσταση ανά πηγή δυσκολίας. «Καλύτερος» σημαίνει τη μικρότερη κύρια μετρική
στο $d=10$ χωρίς μεροληψία· στο διπλό πηγάδι ο ULA έχει μικρότερο $b^2$, αλλά λόγω
μεροληψίας.}
\label{tab:recommendation}
\end{table}

\section{Επιλογή Κατάλληλου Αλγορίθμου}
\label{sec:concl-when}

\paragraph{ULA --- μόνο ως βάση αναφοράς.}
Ο ULA είναι ο απλούστερος δειγματολήπτης Langevin, αλλά έχει μόνιμη μεροληψία σε κάθε
κατανομή, και μεγαλύτερη αλυσίδα δεν τη μειώνει. Είναι τελευταίος σε έξι από τις οκτώ
κατανομές στο $d=10$. Στις κατανομές όπου η μεροληψία του είναι κάτω από το κατώφλι του
$b^2$, φαίνεται στις ουρές. Είναι χρήσιμη βάση αναφοράς, αλλά η παραγωγική χρήση θα
πρέπει να επιλέγει έναν από τους άλλους δύο.

\paragraph{MALA --- σωστός, αλλά χρειάζεται μεγάλη αλυσίδα σε υψηλή διάσταση.}
Ο MALA δεν έχει μεροληψία και συγκλίνει με τον ιδανικό ρυθμό σε τέσσερις από τις
επτά κατανομές της μελέτης σύγκλισης. Το βήμα του όμως πρέπει να
μικραίνει με τη διάσταση, οπότε σε σταθερό μήκος αλυσίδας χειροτερεύει, και στο
$d=1000$ το ESS ανά δευτερόλεπτο του είναι το μικρότερο από τους τρεις σε έξι από τις
επτά κατανομές όπου έχει έγκυρη διαμόρφωση. Είναι η ασφαλής
επιλογή όταν η απουσία μεροληψίας μετρά περισσότερο από το κόστος.

\paragraph{BAOAB --- η προεπιλογή.}
Στο $d=10$ ο BAOAB έχει τη μικρότερη KS σε επτά από τις οκτώ κατανομές και το
μεγαλύτερο ESS σε πέντε (Πίνακας~\ref{tab:best-all}). Η απόδοσή του μένει σταθερή έως
$d=1000$ στη Γκαουσιανή, στην ανισοτροπική Γκαουσιανή, στο μείγμα και στο διπλό πηγάδι. Το πλεονέκτημά του φαίνεται και στο
ESS ανά δευτερόλεπτο: στη Γκαουσιανή είναι περίπου $3$, $12$ και $41$ φορές πιο
αποδοτικός από τον MALA στο $d=10$, $100$ και $1000$ (Ενότητα~\ref{sec:exec-time}). Ο
λόγος είναι η ορμή: με μικρή τριβή η αλυσίδα συνεχίζει στην ίδια κατεύθυνση για
μερικά βήματα και διασχίζει γρήγορα τις φαρδιές κατευθύνσεις. Η τριβή πρέπει να είναι
μικρή ($\gamma = 0.5$) σε έξι από τις οκτώ κατανομές· στο διπλό πηγάδι μεγαλύτερη
τριβή ($\gamma = 2$) είναι καλύτερη.

\section{Κατευθύνσεις για Μελλοντική Έρευνα}
\label{sec:concl-future}

\paragraph{Προρρύθμιση.}
Και οι τρεις αλγόριθμοι επιδέχονται προρρύθμιση που προσαρμόζεται στην τοπική
καμπυλότητα του $U$. Στην ανισοτροπική Γκαουσιανή ένας καλά επιλεγμένος προρρυθμιστής
αντισταθμίζει τη διαφορά κλίμακας μεταξύ συντεταγμένων. Στο χωνί και στην Cauchy ένας
προρρυθμιστής που εξαρτάται από τη θέση θα μπορούσε να μετριάσει το πρόβλημα του
μικρού ευσταθούς βήματος και της κλίσης που σβήνει. Ο Riemann manifold MALA
\citep{girolami2011riemann} και ο προρρυθμισμένος underdamped Langevin θα ήταν φυσικές
γραμμές βάσης σε μια τέτοια μελέτη.

\paragraph{Μεγαλύτερες αλυσίδες σε υψηλή διάσταση.}
Η μελέτη σύγκλισης έγινε στο $d=20$. Σε μεγαλύτερη διάσταση το βήμα του MALA είναι
μικρότερο, οπότε δεν ξέρουμε πόσο μεγάλη αλυσίδα χρειάζεται για να περάσει το
κατώφλι, ούτε αν η αποτυχία όλων των αλγορίθμων στην Cauchy και στο χωνί για μεγάλο $d$
αλλάζει με πολύ μεγαλύτερες αλυσίδες.

\paragraph{Σύγκριση με τη Hamiltonian Monte Carlo.}
Η Hamiltonian Monte Carlo \citep{neal2011mcmc} συνδυάζει την ορμή του BAOAB με έναν
έλεγχο Metropolis. Τα αποτελέσματα της εργασίας --- ο BAOAB κερδίζει μέσω της ορμής, ο
MALA μέσω της απουσίας μεροληψίας --- κάνουν φυσική την ερώτηση αν ένας τέτοιος
συνδυασμός κρατά και τα δύο πλεονεκτήματα στις ίδιες κατανομές.

\paragraph{Προσαρμοστική τριβή στον BAOAB.}
Η μικρή τριβή είναι καλύτερη στις περισσότερες κατανομές, αλλά όχι στο διπλό πηγάδι.
Ένα σχήμα που προσαρμόζει το $\gamma$ κατά τη διάρκεια της εκτέλεσης θα μπορούσε να
αποφύγει τη ρύθμιση ανά κατανομή. Αν ένα τέτοιο σχήμα μπορεί να ταιριάξει ή να υπερβεί
την απόδοση σταθερού $\gamma$ είναι ανοιχτό ερώτημα.
